\bibitem{refKuchuganov}
Кучуганов, А. В. Методология семантического анализа и поиска графической информации : автореферат диссертации на соискание ученой степени доктора технических наук : специальность 05.13.01 «Системный анализ, управление и обработка информации (в промышленности)» / А. В. Кучуганов. – Ижевск, 2018. – Текст : непосредственный.
\bibitem{refDavletov}
Давлетов, А. Р. Современные методы машинного обучения и технология OCR для автоматизации обработки документов / А. Р. Давлетов. – Текст : электронный // Вестник науки. – 2023. – № 10 (67). – URL: https://cyberleninka.ru/article/n/sovremennye-metody-mashinnogo-obucheniya-i-tehnologiya-ocr-dlya-avtomatizatsii-obrabotki-dokumentov (дата обращения: 22.05.2026).
\bibitem{refRakov}
Раков, Н. С. Машинное зрение / Н. С. Раков, С. В. Пальмов. – Текст : электронный // Форум молодых ученых. – 2018. – № 4 (20). – URL: https://cyberleninka.ru/article/n/mashinnoe-zrenie (дата обращения: 22.05.2026).
\bibitem{refRublevsky}
Рублевский, Р. Ландшафт основных терминов в области генеративного AI, их взаимосвязь и употребление / Р. Рублевский. – Текст : электронный // Хабр. – 2025. – 25 сент. – URL: https://habr.com/ru/articles/950600/ (дата обращения: 22.05.2026).
\bibitem{refAkhil}
Akhil, S. Overview of Tesseract OCR engine / S. Akhil. – Текст : электронный // ResearchGate. – 2016. – Dec. – URL: https://www.researchgate.net/publication/315834331 (дата обращения: 22.05.2026).
\bibitem{refLayoutParser}
Shen, Z. LayoutParser: A Unified Toolkit for Deep Learning Based Document Image Analysis / Z. Shen [et al.]. – Текст : электронный // arXiv.org. – 2021. – arXiv:2103.15348. – URL: https://arxiv.org/abs/2103.15348 (дата обращения: 22.05.2026).
\bibitem{refLayoutLM}
Xu, Y. LayoutLM: Pre-training of Text and Layout for Document Image Understanding / Y. Xu, M. Li, L. Cui [et al.]. – Текст : электронный // arXiv.org. – 2019. – arXiv:1912.13318. – URL: https://arxiv.org/abs/1912.13318 (дата обращения: 22.05.2026).
\bibitem{refSmith}
Smith, R. An Overview of the Tesseract OCR Engine / R. Smith // Ninth International Conference on Document Analysis and Recognition (ICDAR 2007). – Washington : IEEE Computer Society, 2007. – Vol. 2. – P. 629--633. – DOI: 10.1109/ICDAR.2007.4376991. – URL: https://doi.org/10.1109/ICDAR.2007.4376991 (дата обращения: 22.05.2026).
\bibitem{refTrinajstic}
Trinajstić, N. Chemical Graph Theory / N. Trinajstić. – 2nd ed. – Boca Raton : CRC Press, 1992. – 364 p.

\bibitem{refWarr}
Warr, W. A. Representation of Chemical Structures / W. A. Warr // Wiley Interdisciplinary Reviews: Computational Molecular Science. – 2011. – Vol. 1, no. 4. – P. 557--579. – DOI: 10.1002/wcms.36.

\bibitem{refWeininger1}
Weininger, D. SMILES, a Chemical Language and Information System. 1. Introduction to Methodology and Encoding Rules / D. Weininger // Journal of Chemical Information and Computer Sciences. – 1988. – Vol. 28, no. 1. – P. 31--36. – DOI: 10.1021/ci00057a005.

\bibitem{refWeininger2}
Weininger, D. SMILES. 2. Algorithm for Generation of Unique SMILES Notation / D. Weininger, A. Weininger, J. L. Weininger // Journal of Chemical Information and Computer Sciences. – 1989. – Vol. 29, no. 2. – P. 97--101. – DOI: 10.1021/ci00062a008.

\bibitem{refDalby}
Dalby, A. Description of Several Chemical Structure File Formats Used by Computer Programs Developed at Molecular Design Limited / A. Dalby, J. G. Nourse, W. D. Hounshell [et al.] // Journal of Chemical Information and Computer Sciences. – 1992. – Vol. 32, no. 3. – P. 244--255. – DOI: 10.1021/ci00007a012.

\bibitem{refInChI}
Heller, S. R. InChI, the IUPAC International Chemical Identifier / S. R. Heller, A. McNaught, I. Pletnev [et al.] // Journal of Cheminformatics. – 2015. – Vol. 7. – Art. 23. – DOI: 10.1186/s13321-015-0068-4.

\bibitem{refOBoyleCanonical}
O'Boyle, N. M. Towards a Universal SMILES Representation – A Standard Method to Generate Canonical SMILES Based on the InChI / N. M. O'Boyle // Journal of Cheminformatics. – 2012. – Vol. 4. – Art. 22. – DOI: 10.1186/1758-2946-4-22.

\bibitem{refWikiOCSR}
Optical chemical structure recognition. – Текст : электронный // Wikipedia. – URL: $https://en.wikipedia.org/wiki/Optical_chemical_structure_recognition$ (дата обращения: 22.05.2026).

\bibitem{refUML}
OMG Unified Modeling Language (OMG UML), Version 2.5.1 / Object Management Group. – 2017. – URL: https://www.omg.org/spec/UML/2.5.1/PDF (дата обращения: 22.05.2026).

\bibitem{refMolScribe}
MolScribe: Robust Molecular Structure Recognition with Image-To-Graph Generation / Y. Qian, J. Guo, Z. Tu [et al.]. – Текст : электронный // arXiv.org. – 2022. – arXiv:2205.14311. – URL: https://arxiv.org/pdf/2205.14311 (дата обращения: 22.05.2026).

\bibitem{refDECIMER}
Rajan, K. DECIMER 1.0: deep learning for chemical image recognition using transformers / K. Rajan [et al.]. – Текст : электронный // ResearchGate. – 2021. – Aug. – URL: $https://www.researchgate.net/publication/353957394_DECIMER_10_deep_learning_for_chemical_image_recognition_using_transformers$ (дата обращения: 22.05.2026).

\bibitem{refDECIMERJournal}
Rajan, K. DECIMER 1.0: deep learning for chemical image recognition using transformers / K. Rajan, A. Zielesny, C. Steinbeck. – Текст : электронный // Journal of Cheminformatics. – 2021. – Vol. 13. – DOI: 10.1186/s13321-021-00538-8. – URL: https://doi.org/10.1186/s13321-021-00538-8 (дата обращения: 22.05.2026).

\bibitem{refQwen3VL}
Qwen Team. Qwen3-VL Technical Report / Qwen Team. – Текст : электронный // arXiv.org. – 2025. – arXiv:2511.21631. – URL: https://arxiv.org/abs/2511.21631 (дата обращения: 22.05.2026).

\bibitem{refCAMP}
Fifty, C. In-Context Learning for Few-Shot Molecular Property Prediction / C. Fifty, J. Leskovec, S. Thrun. – Текст : электронный // arXiv.org. – 2023. – arXiv:2310.08863. – URL: https://arxiv.org/abs/2310.08863 (дата обращения: 22.05.2026).

\bibitem{refGoF}
Gamma, E. Design Patterns: Elements of Reusable Object-Oriented Software / E. Gamma, R. Helm, R. Johnson, J. Vlissides. – Reading : Addison-Wesley, 1995. – 395 p. – ISBN 0-201-63361-2. – Текст : непосредственный.

\bibitem{refPEP8}
van Rossum, G. PEP 8 -- Style Guide for Python Code / G. van Rossum, B. Warsaw, A. Coghlan. – Текст : электронный // Python.org. – URL: https://peps.python.org/pep-0008/ (дата обращения: 23.05.2026).

\bibitem{refLineAngle}
Chapter 1 - Organic Chemistry Review / Hydrocarbons. Section 1.4 Condensed Structural and Line-Angle Formulas. – Текст : электронный // Hostos Community College Library Guides. – URL: https://guides.hostos.cuny.edu/che120 (дата обращения: 22.05.2026).

\bibitem{refRLSA}
Wong, K. Y. Document Analysis System / K. Y. Wong, R. G. Casey, F. M. Wahl // IBM Journal of Research and Development. – 1982. – Vol. 26, no. 6. – P. 647--656. – DOI: 10.1147/rd.266.0647.

\bibitem{refXYCut}
Ha, J. Recursive X-Y Cut Using Bounding Boxes of Connected Components / J. Ha, R. M. Haralick, I. T. Phillips // Proceedings of the Third International Conference on Document Analysis and Recognition. – 1995. – Vol. 2. – P. 952--955. – DOI: 10.1109/ICDAR.1995.602058.

\bibitem{refStarovoytov}
Старовойтов, В. В. Получение и обработка изображений на ЭВМ : учебно-методическое пособие / В. В. Старовойтов, Ю. И. Голуб. – Минск : БНТУ, 2018. – 204 с. – ISBN 978-985-550-770-4. – Текст : непосредственный.

\bibitem{refOCRMetrics}
Smith, R. An Overview of the Tesseract OCR Engine / R. Smith // Ninth International Conference on Document Analysis and Recognition (ICDAR 2007). – Washington : IEEE Computer Society, 2007. – Vol. 2. – P. 629–633. – DOI: 10.1109/ICDAR.2007.4376991. – URL: https://doi.org/10.1109/ICDAR.2007.4376991 (дата обращения: 22.05.2026).

\bibitem{refLlava}
Liu, H. Visual Instruction Tuning / H. Liu, C. Li, Q. Wu, Y. J. Lee. – Текст : электронный // arXiv.org. – 2023. – arXiv:2305.11845. – URL: https://arxiv.org/abs/2305.11845 (дата обращения: 22.05.2026).
